\section{Introduction and precise statements}\label{sec:introduction}

Planar critical models are expected to have conformally invariant scaling
limits.  For random-cluster models this prediction identifies the law of an
interface, the joint law of all surrounding loops, and crossing
probabilities in arbitrary shapes.  These are distinct conclusions: a
limit for one interface does not by itself identify its coupling with
other interfaces, and convergence of loop traces on compact sets does not
by itself retain multiplicities, nesting, or complete traversals.

Schramm's characterization of conformally invariant interfaces led to
Schramm--Loewner evolution \citep{Schramm2000}.  Cardy's prediction for
critical percolation crossings \citep{Cardy1992} was proved by Smirnov for
site percolation on the triangular lattice \citep{Smirnov2001}.  The
random-cluster Ising model has a discrete holomorphic observable
\citep{Smirnov2010}, and its interface convergence is established in
\citet{ChelkakDCHonglerKemppainenSmirnov2014}.  Conformal loop ensembles
provide the corresponding continuum laws for complete families of
interfaces \citep{SheffieldWerner2012,Sheffield2009Exploration}.
The results below concern the nearest-neighbor square lattice, with its
ordinary Euclidean embedding and its original random-cluster probability
law.

\subsection{The model and its boundary conventions}
\label{intro:model}

Fix a real number $q\in[1,4]$ and put
\begin{equation}\label{intro:parameters}
 p_q=\frac{\sqrt q}{1+\sqrt q},
 \qquad
 \kappa(q)=\frac{4\pi}{\arccos(-\sqrt q/2)}.
\end{equation}
The inverse cosine has values in $[0,\pi]$.  Thus $4\le\kappa(q)\le6$,
with $\kappa(1)=6$ and $\kappa(4)=4$.  On a finite graph $G=(V,E)$,
let $\xi$ be a partition of designated boundary vertices.  For
$A\subseteq E$, let $k_\xi(A)$ be the number of components of $(V,A)$
after identification within each block of $\xi$, including isolated
vertices.  The random-cluster law is
\begin{equation}\label{intro:finite-law}
 \phi_{G,q}^{\xi}(A)
 =\frac{1}{Z_{G,q}^{\xi}}
 p_q^{|A|}(1-p_q)^{|E|-|A|}q^{k_\xi(A)}.
\end{equation}
An \emph{actual open path} always consists of open edges of $G$.
Identification in $\xi$ is used to compute the probability weight and
does not supply an additional edge in such a path.

Let $D$ be a bounded Jordan domain.  A lattice approximation at mesh
$\delta_n\downarrow0$ is the bounded interior $D_n$ of a simple closed
nearest-neighbor polygon on $\delta_n\Z^2$.  Its graph $G_n$ has every
lattice vertex in $\overline{D_n}$ and every nearest-neighbor edge whose
closed segment is contained in $\overline{D_n}$.  In particular, boundary
edges are part of $G_n$.

For two marked boundary points $a,b$, we require counterclockwise
parametrizations $\beta_n,\beta:[0,1]\to\C$, injective on $[0,1)$,
such that
\begin{gather}\label{intro:marked-jordan}
 \beta_n(0)=\beta_n(1)=a_n,\qquad \beta_n(1/2)=b_n,\qquad
 \beta(0)=\beta(1)=a,\qquad \beta(1/2)=b,\\
 \sup_{t\in[0,1]}|\beta_n(t)-\beta(t)|\longrightarrow0.\notag
\end{gather}
The free arc $F_n$ goes counterclockwise from $a_n$ to $b_n$; the other
arc $W_n$ goes counterclockwise from $b_n$ to $a_n$.  Wire together all
vertices of $W_n$, including the two marked endpoints.  All other
boundary vertices remain singleton blocks.  Remove from the sampled edge
set the boundary edges of $W_n$ and declare those edges open for the
interface drawing.  This is equivalent on all remaining edges to
\eqref{intro:finite-law} with the same wired partition: the endpoints of
each removed edge were already identified.  Edges on $F_n$ remain random.

The medial graph has vertices at the midpoints of primal edges, with
medial edges joining the midpoints of consecutive sides of a lattice
square.  At an open primal edge the four incident medial half-edges are
paired in quarter turns around the adjacent dual faces; at a closed
primal edge they are paired around the primal endpoints.  At the
boundary one uses the turns between the primal wired arc and its
complementary dual wired arc.  There is a unique medial strand between
the two changes of boundary condition.  We orient it from $a_n$ to $b_n$
and denote it by $\eta_n$.  This is the usual Dobrushin construction
\citep[Section~5.2.2]{DC2017Lectures}.  Occurrence order at a twice-visited
medial vertex is part of the curve.  The tile drawings used later preserve
this order and differ from the straight medial drawing by at most a
constant times $\delta_n$ in uniform curve distance.

For continuous oriented curves $\gamma_1,\gamma_2:[0,1]\to\C$, write
\begin{equation}\label{intro:curve-metric}
 d_{\mathrm{curv}}(\gamma_1,\gamma_2)
 =\inf_{\varphi_1,\varphi_2}
   \sup_{0\le t\le1}
   |\gamma_1(\varphi_1(t))-\gamma_2(\varphi_2(t))|,
\end{equation}
where the infimum is over increasing homeomorphisms of $[0,1]$.
Zero-distance curves are identified.  Chordal $\SLE_\kappa$ in $(D;a,b)$
is the conformal image of the upper-half-plane Loewner trace driven by
$\sqrt\kappa B_t$ from $0$ to $\infty$, with $B$ a standard Brownian
motion.  The remaining positive dilation of the conformal coordinate
does not change this law.

\subsection{Complete loop collections}
\label{intro:collections}

Let $\phi_q^0$ be the weak limit of \eqref{intro:finite-law} in increasing
square boxes of $\Z^2$ with free boundary conditions.  Form its complete
medial interface collection, using the same local pairings as above, and
then scale the embedding by $\delta$.  Denote the resulting collection
by $\Gamma_\delta$.  It contains every loop at every nesting level.
The elementary circuit argument in Section~\ref{sec:preliminaries}
shows that its discrete interfaces are almost surely finite, also at
$q=4$.

Individual loops are unrooted and unoriented continuous maps of the
circle into the Riemann sphere $\widehat\C$.  Their distance is the
infimum of the uniform spherical distance over circular
orientation-preserving homeomorphisms and reversal, with zero-distance
representatives identified.  We retain distinct loops even when their
images coincide.  For two collections, an $\eps$-matching is a partial
bijection that covers all loops of spherical diameter greater than
$\eps$ in both collections, has matched-loop distance at most $\eps$,
and leaves only loops of diameter at most $\eps$ unmatched.  Convergence
of collections means the existence of such matchings with
$\eps\downarrow0$.  The continuum law in the next statement is the
standard nested whole-plane $\CLE_\kappa$ in this topology.

\begin{theorem}[Interfaces and all nested loops]
\label{intro:main}
For every fixed real $1\le q\le4$, the following assertions hold.
\begin{enumerate}[label=\textup{(\roman*)}]
\item For every bounded Jordan domain $D$ with distinct boundary marks
$a,b$, and every approximation satisfying
\eqref{intro:marked-jordan}, the Dobrushin strand $\eta_n$ converges in
law, in \eqref{intro:curve-metric}, to chordal $\SLE_{\kappa(q)}$ from
$a$ to $b$ in $D$.
\item Under the specified free infinite-volume law,
$\Gamma_\delta$ converges in law as $\delta\downarrow0$ to the nested
whole-plane $\CLE_{\kappa(q)}$, in the spherical collection topology
just defined.  The limit is M\"obius invariant.
\end{enumerate}
Both conclusions hold along every mesh sequence.  No uniformity in $q$
is asserted.
\end{theorem}

Theorem~\ref{intro:main} establishes the SLE/CLE conformal-invariance
prediction for critical random-cluster models on the square lattice
throughout $q\in[1,4]$.  The word ``critical'' refers to the edge
parameter $p_q$.  At $q=4$ the limits are chordal $\SLE_4$ and the
standard simple nested whole-plane $\CLE_4$.

The endpoint proof uses a whole-plane height-field theorem, with its
precise law and symmetry input verified in
Section~\ref{sec:q4-rotation-repair}.  That theorem supplies no
finite-domain field convergence, conditional convergence after an
exploration, or interface convergence.  Sections~\ref{sec:q4regular}--
\ref{sec:q4identification} prove these further implications.  The
strict angular and six-arm estimates used for $q<4$ are not extended to
$q=4$ by a limiting argument.

\subsection{Square-lattice bond crossings}
\label{intro:crossings}

For four distinct counterclockwise boundary marks $a,b,c,d$ of $D$,
require the same uniform convergence as in
\eqref{intro:marked-jordan}, with the respective four marks at
parameters $0,1/4,1/2,3/4$.  Sample \emph{all} edges of $G_n$
independently with probability $1/2$, including every boundary edge, and
make no identifications or external connections.  Let $C_n$ be the
event of an actual open path from the closed counterclockwise arc
$[a_n,b_n]$ to the closed counterclockwise arc $[c_n,d_n]$.

\begin{theorem}[Cardy's formula for square-lattice bond percolation]
\label{intro:cardy}
Let $\psi:D\to\H$ send $a,c,d$ to $0,1,\infty$, and put
$x=\psi(b)\in(0,1)$.  Under the boundary and approximation conventions
above,
\begin{equation}\label{intro:cardy-formula}
 \lim_{n\to\infty}\mathbb P(C_n)
 =\frac{\displaystyle\int_0^x[t(1-t)]^{-2/3}\,\dd t}
        {\displaystyle\int_0^1[t(1-t)]^{-2/3}\,\dd t}.
\end{equation}
\end{theorem}

Theorem~\ref{intro:cardy} proves the square-lattice bond-percolation form
of Cardy's prediction with the free-boundary convention specified here.
We obtain it directly from the alternating four-change calculation and
a deterministic comparison on a common lattice graph.  In particular,
its proof supplies the boundary transfer for the actual crossing event;
it does not infer continuity of that event from a curve limit.

\subsection{Proof structure and technical contributions}

For $1\le q<4$, use the parameters
\begin{equation}\label{intro:lambda-rho}
 d=\sqrt q=2\cos\lambda,\qquad
 \rho=\lambda/\pi,\qquad 0<\lambda\le\pi/3.
\end{equation}
The first part of the proof is a four-change boundary calculation.  In
the upper half-plane, put the boundary changes at $0,\chi,1,\infty$.
For the open cap normalization, in which the four unmatched ends carry
no additional pairing weights, the connection probability is
\begin{equation}\label{intro:fourchange-formula}
 f(\chi)=\frac{I_C(\chi)}{I_B(\chi)+I_C(\chi)},
 \quad
 \begin{cases}
 I_B(\chi)=\displaystyle\int_\chi^1
 u^{\rho-1}|u-\chi|^{\rho-1}|u-1|^{1-3\rho}\,\dd u,\\[3pt]
 I_C(\chi)=\displaystyle\int_1^\infty
 u^{\rho-1}|u-\chi|^{\rho-1}|u-1|^{1-3\rho}\,\dd u.
 \end{cases}
\end{equation}
Separate deterministic cap wires on the two designated intervals change
this to $f/(f+d(1-f))$.  Keeping these two normalizations distinct is
essential when $d\ne1$.

\paragraph{Positive comparisons before complex observables.}
Section~\ref{sec:preliminaries} identifies the capped tile law exactly
with the random-cluster law and establishes the collar and passage-count
comparisons used later.  It keeps actual components distinct from
boundary partition blocks.  Bulk box crossing at $q\le4$ and the stronger
up-to-boundary crossing estimates for $q<4$ have different roles
\citep{DCLM2019Isoradial,DMT2021Fractal}.  Section~\ref{sec:exponents}
derives the local exponent inequalities needed to control corner, tip,
and source factors.  All complex currents are subsequently estimated in
these positive laws.

\paragraph{Angular transfer and boundary holomorphicity.}
Sections~\ref{sec:angular} and \ref{sec:current} construct the transfer
operators, their reflected norms, and the path-current primitive.
The argument distinguishes a common complex continuation radius from
norm constants that may depend on a vector in the cyclic core.  A
scalar spectral factorization then suppresses the angular modes that
would obstruct holomorphicity.  The boundary calculation in
Section~\ref{sec:fourchange} is initially conditional on this
holomorphicity.  Sections~\ref{sec:pairedwall}--\ref{sec:holomorphic}
discharge the condition by a distant-probe comparison, closing finitely
many wall gaps, and a folded mixed-stencil estimate.  The folds retain
the external states in their reflected norms; local exponent savings
pay for extraction and normalization costs.  This produces
\eqref{intro:fourchange-formula} without assuming a lattice conformal
invariance theorem.

\paragraph{Moving domains and the chordal driver.}
The fixed-shape formula must remain useful after a random exploration.
Section~\ref{sec:movingtests} proves this transfer using a finite chart
approximation and an averaged six-arm estimate under the original
switch law.  It does not require a uniform six-arm theorem in every
possible explored domain.  Section~\ref{sec:chordal} first normalizes a
shrinking test into a bounded martingale and then transfers it to the
true Dobrushin law.  Two tests identify the driver and its bracket.
Tightness and a separate traversal argument give the Euclidean uniform
curve conclusion of Theorem~\ref{intro:main}(i) for $q<4$.

\paragraph{Joint branches and loop identities.}
Sections~\ref{sec:peeling}--\ref{sec:loops} construct a target tree by
uniform-walled peeling.  They establish the reflected angular dynamics,
exclude forcepoint variation supported on endpoint times, and identify
independent continuation noises after separation, conditional on the
common causal history at separation.  The resulting branching SLE description is then matched
to the original discrete loops: winding information determines loop
labels, while additional estimates determine full traversal order.
An escape bound upgrades convergence in bounded windows to the
spherical matching topology for every nesting generation.

\paragraph{The endpoint field argument.}
At $q=4$ the loop height uses independent fair signs and unit jumps.
Section~\ref{sec:q4-rotation-repair} extracts a bijection of selected
loops from the public FK homotopy coupling
\citep{DCKKMO2026Rotation}.  Even sign moments depend only on loops
that separate at least two distinct pairs of endpoints.  These loops
are macroscopic, and the bijection preserves their winding differences
up to one harmless sign per loop.  This supplies the endpoint symmetry
input for the whole-plane six-vertex field theorem
\citep{SixVertexGFF2026}, including the multiplicities that a loop-set
distance alone would not determine.
Section~\ref{sec:q4tools} combines binary transport and reflection
positivity with the whole-plane height input to control double
contacts.  The error estimates are relative to the square of the
four-arm probability; an absolute error would be insufficient.
Section~\ref{sec:q4regular} derives the Dirichlet field in regular
domains by translating separated charge groups and determining the
harmonic boundary term.  Section~\ref{sec:q4rough} then treats rough
explored pockets by a conditional moment comparison, with contour
injections at each numerical nesting level and a regular-domain
sandwich.  The final identification verifies locality, thinness,
bounded harmonic values, and the relevant topology before invoking
two-valued local-set uniqueness
\citep{AruSepulvedaWerner2019,AruLupuSepulveda2020}.

\subsection{Conventions on constants and limits}

Every theorem is for a fixed $q$.  Constants may depend on $q$,
macroscopic shapes, positive separations, and an explicitly fixed
finite set of tests.  Mesh limits are taken with those data fixed;
auxiliary separations, gaps, cutoffs, and localization parameters are
removed in the stated order.  A claim uniform over random exploration
histories is made only where the corresponding comparison has been
proved.  In particular, convergence in probability over histories is
not replaced by a deterministic estimate for every history.

The proof uses established public results at their stated hypotheses.
It uses no theorem about an annealed random planar map, no random
embedding, and no transfer from a random-map metric limit to the square
lattice.  The finite Dobrushin experiments and the free infinite-volume
experiment are different probability spaces; the theorems assert their
separate convergence and require no coupling between their samples.
